\section{周期性微扰}

\subsection{单色驱动}

若外场为
\begin{equation}
  F(\mathbf{r},t)
  =F_{\mathbf{q}}\,e^{\mathrm{i}(\mathbf{q}\cdot\mathbf{r}-\omega t)}
  +\mathrm{c.c.},
\end{equation}
则在稳态下诱导期望值以同频振荡（线性级）：
\begin{equation}
  \delta\langle B\rangle_{\mathbf{q},\omega}
  =\chi_{BA}(\mathbf{q},\omega)\,F_{\mathbf{q}}.
\end{equation}
实验上的光学吸收、介电测量、非弹性散射等，往往直接探测特定 $(\mathbf{q},\omega)$ 处的响应。

\subsection{绝热打开}

为确保因果与绝热切换，常令外场含因子 $e^{\eta t}$（$\eta\to0^+$）从 $t=-\infty$ 缓慢打开。这等价于在频率中取 $\omega\to\omega+\mathrm{i}\eta$，从而挑选 retarded 函数。

\subsection{耗散}

功率吸收与响应虚部相关。例如，对电偶极驱动，吸收正比于电导率实部，而电导率又由电流--电流响应给出。这与涨落耗散定理一致：虚部编码耗散，也编码涨落谱。
